
\documentclass{article}
\usepackage{colortbl}
\usepackage{makecell}
\usepackage{multirow}
\usepackage{supertabular}

\begin{document}

\newcounter{utterance}

\centering \large Interaction Transcript for game `dond', experiment `semi', episode 19 with Qwen3.6{-}35B{-}A3B{-}FP8{-}without{-}reasoning.
\vspace{24pt}

{ \footnotesize  \setcounter{utterance}{1}
\setlength{\tabcolsep}{0pt}
\begin{supertabular}{c@{$\;$}|p{.15\linewidth}@{}p{.15\linewidth}p{.15\linewidth}p{.15\linewidth}p{.15\linewidth}p{.15\linewidth}}
    \# & $\;$A & \multicolumn{4}{c}{Game Master} & $\;\:$B\\
    \hline

    \theutterance \stepcounter{utterance}  
    & & \multicolumn{4}{p{0.6\linewidth}}{
        \cellcolor[rgb]{0.9,0.9,0.9}{
            \makecell[{{p{\linewidth}}}]{
                \texttt{\tiny{[P1$\langle$GM]}}
                \texttt{Παίζετε ένα συνεργατικό παιχνίδι διαπραγμάτευσης στο οποίο πρέπει να συμφωνήσετε με έναν άλλο παίκτη σχετικά με τον τρόπο διαίρεσης ενός συνόλου αντικειμένων μεταξύ σας.} \\
\\ 
\texttt{Κανόνες:} \\
\texttt{(a) Εσείς και ο άλλος παίκτης λαμβάνετε ένα σύνολο αντικειμένων. Σας δίνεται επίσης μια μυστική συνάρτηση τιμής που αντιπροσωπεύει πόσο εκτιμάτε κάθε τύπο αντικειμένου.} \\
\texttt{(b) Ανταλλάσσετε μηνύματα με τον άλλο παίκτη για να συμφωνήσετε σχετικά με το ποιος παίρνει ποια αντικείμενα. Μπορείτε να στείλετε το πολύ 5 μηνύματα ο καθένας ή να τερματίσετε νωρίς κάνοντας μια μυστική πρόταση ανά πάσα στιγμή.} \\
\texttt{(c) Σας ζητείται να υποβάλετε μια μυστική πρόταση που υποδεικνύει τα αντικείμενα που θέλετε, μορφοποιημένη σε τετράγωνες αγκύλες ως εξής: "{[}Πρόταση: <αριθμός> <όνομα αντικειμένου>, <αριθμός> <όνομα αντικειμένου>, <...>{]}"} \\
\texttt{(d) Εάν οι προτάσεις σας είναι συμπληρωματικές, δηλαδή υπάρχουν αρκετά αντικείμενα για να ικανοποιήσουν και τις δύο προτάσεις, κάθε παίκτης λαμβάνει βαθμολογία με βάση το άθροισμα των τιμών των αντικειμένων που έλαβε. Διαφορετικά, και οι δύο παίρνετε μηδέν πόντους.} \\
\texttt{(e) Ο στόχος σας είναι να μεγιστοποιήσετε το άθροισμα της βαθμολογίας σας και της βαθμολογίας του άλλου παίκτη.} \\
\\ 
\texttt{Ας ξεκινήσουμε.} \\
\\ 
\texttt{Το σύνολο των διαθέσιμων αντικειμένων είναι:} \\
\\ 
\texttt{2 λάμπες, 3 λουλούδια, 1 μπανάνα} \\
\\ 
\texttt{Η μυστική συνάρτηση τιμής σας είναι:} \\
\\ 
\texttt{λάμπα: 4, λουλούδι: 0, μπανάνα: 2} \\
\\ 
\texttt{ΣΗΜΑΝΤΙΚΟ: Τα μηνύματά σας, εκτός εάν είναι μυστική πρόταση, μεταδίδονται απευθείας στον άλλο παίκτη, επομένως μην συμπεριλάβετε καμία απάντηση στους κανόνες ή κείμενο που ανακοινώνει το μήνυμά σας. Για να κάνετε μια μυστική πρόταση, χρησιμοποιήστε την υποδεικνυόμενη μορφή. Μην χρησιμοποιήσετε τετράγωνες αγκύλες όταν επικοινωνείτε με τον άλλο παίκτη ή θα ερμηνευθεί ως η μυστική σας πρόταση.} \\
            }
        }
    }
    & \\ \\

    \theutterance \stepcounter{utterance}  
    & \multicolumn{4}{p{0.6\linewidth}}{
        \cellcolor[rgb]{0.9,0.9,0.9}{
            \makecell[{{p{\linewidth}}}]{
                \texttt{\tiny{[P1$\rangle$GM]}}
                \texttt{Γεια σου! Για να μεγιστοποιήσουμε το συνολικό μας κέρδος, ας δούμε πρώτα τις αξίες μας. Εγώ εκτιμώ τις λάμπες σε 4 μονάδες και τη μπανάνα σε 2 μονάδες, ενώ τα λουλούδια για μένα δεν έχουν καμία αξία (0).} \\
\\ 
\texttt{Επειδή δεν γνωρίζω τις προτιμήσεις σου, θα ήθελα να μάθω τι αντικείμενα εκτιμάς περισσότερο. Αν μπορείς, πες μου ποια αντικείμενα θέλεις εσύ, ώστε να δούμε πώς μπορούμε να τα μοιραστούμε χωρίς σύγκρουση.} \\
\\ 
\texttt{Ποια αντικείμενα θα ήθελες εσύ;} \\
            }
        }
    }
    & & \\ \\

    \theutterance \stepcounter{utterance}  
    & & & \multicolumn{2}{p{0.3\linewidth}}{
        \cellcolor[rgb]{0.9,0.9,0.9}{
            \makecell[{{p{\linewidth}}}]{
                \texttt{\tiny{[GM$|$GM]}}
                \texttt{continue} \\
            }
        }
    }
    & & \\ \\

    \theutterance \stepcounter{utterance}  
    & & & \multicolumn{2}{p{0.3\linewidth}}{
        \cellcolor[rgb]{0.9,0.9,0.9}{
            \makecell[{{p{\linewidth}}}]{
                \texttt{\tiny{[GM$|$GM]}}
                \texttt{Γεια σου! Για να μεγιστοποιήσουμε το συνολικό μας κέρδος, ας δούμε πρώτα τις αξίες μας. Εγώ εκτιμώ τις λάμπες σε 4 μονάδες και τη μπανάνα σε 2 μονάδες, ενώ τα λουλούδια για μένα δεν έχουν καμία αξία (0).} \\
\\ 
\texttt{Επειδή δεν γνωρίζω τις προτιμήσεις σου, θα ήθελα να μάθω τι αντικείμενα εκτιμάς περισσότερο. Αν μπορείς, πες μου ποια αντικείμενα θέλεις εσύ, ώστε να δούμε πώς μπορούμε να τα μοιραστούμε χωρίς σύγκρουση.} \\
\\ 
\texttt{Ποια αντικείμενα θα ήθελες εσύ;} \\
            }
        }
    }
    & & \\ \\

    \theutterance \stepcounter{utterance}  
    & & \multicolumn{4}{p{0.6\linewidth}}{
        \cellcolor[rgb]{0.9,0.9,0.9}{
            \makecell[{{p{\linewidth}}}]{
                \texttt{\tiny{[GM$\rangle$P2]}}
                \texttt{Παίζετε ένα συνεργατικό παιχνίδι διαπραγμάτευσης στο οποίο πρέπει να συμφωνήσετε με έναν άλλο παίκτη σχετικά με τον τρόπο διαίρεσης ενός συνόλου αντικειμένων μεταξύ σας.} \\
\\ 
\texttt{Κανόνες:} \\
\texttt{(a) Εσείς και ο άλλος παίκτης λαμβάνετε ένα σύνολο αντικειμένων. Σας δίνεται επίσης μια μυστική συνάρτηση τιμής που αντιπροσωπεύει πόσο εκτιμάτε κάθε τύπο αντικειμένου.} \\
\texttt{(b) Ανταλλάσσετε μηνύματα με τον άλλο παίκτη για να συμφωνήσετε σχετικά με το ποιος παίρνει ποια αντικείμενα. Μπορείτε να στείλετε το πολύ 5 μηνύματα ο καθένας ή να τερματίσετε νωρίς κάνοντας μια μυστική πρόταση ανά πάσα στιγμή.} \\
\texttt{(c) Σας ζητείται να υποβάλετε μια μυστική πρόταση που υποδεικνύει τα αντικείμενα που θέλετε, μορφοποιημένη σε τετράγωνες αγκύλες ως εξής: "{[}Πρόταση: <αριθμός> <όνομα αντικειμένου>, <αριθμός> <όνομα αντικειμένου>, <...>{]}"} \\
\texttt{(d) Εάν οι προτάσεις σας είναι συμπληρωματικές, δηλαδή υπάρχουν αρκετά αντικείμενα για να ικανοποιήσουν και τις δύο προτάσεις, κάθε παίκτης λαμβάνει βαθμολογία με βάση το άθροισμα των τιμών των αντικειμένων που έλαβε. Διαφορετικά, και οι δύο παίρνετε μηδέν πόντους.} \\
\texttt{(e) Ο στόχος σας είναι να μεγιστοποιήσετε το άθροισμα της βαθμολογίας σας και της βαθμολογίας του άλλου παίκτη.} \\
\\ 
\texttt{Ας ξεκινήσουμε.} \\
\\ 
\texttt{Το σύνολο των διαθέσιμων αντικειμένων είναι:} \\
\\ 
\texttt{2 λάμπες, 3 λουλούδια, 1 μπανάνα} \\
\\ 
\texttt{Η μυστική συνάρτηση τιμής σας είναι:} \\
\\ 
\texttt{λάμπα: 0, λουλούδι: 3, μπανάνα: 1} \\
\\ 
\texttt{ΣΗΜΑΝΤΙΚΟ: Τα μηνύματά σας, εκτός εάν είναι μυστική πρόταση, μεταδίδονται απευθείας στον άλλο παίκτη, επομένως μην συμπεριλάβετε καμία απάντηση στους κανόνες ή κείμενο που ανακοινώνει το μήνυμά σας. Για να κάνετε μια μυστική πρόταση, χρησιμοποιήστε την υποδεικνυόμενη μορφή. Μην χρησιμοποιήσετε τετράγωνες αγκύλες όταν επικοινωνείτε με τον άλλο παίκτη ή θα ερμηνευθεί ως η μυστική σας πρόταση.} \\
\\ 
\texttt{Γεια σου! Για να μεγιστοποιήσουμε το συνολικό μας κέρδος, ας δούμε πρώτα τις αξίες μας. Εγώ εκτιμώ τις λάμπες σε 4 μονάδες και τη μπανάνα σε 2 μονάδες, ενώ τα λουλούδια για μένα δεν έχουν καμία αξία (0).} \\
\\ 
\texttt{Επειδή δεν γνωρίζω τις προτιμήσεις σου, θα ήθελα να μάθω τι αντικείμενα εκτιμάς περισσότερο. Αν μπορείς, πες μου ποια αντικείμενα θέλεις εσύ, ώστε να δούμε πώς μπορούμε να τα μοιραστούμε χωρίς σύγκρουση.} \\
\\ 
\texttt{Ποια αντικείμενα θα ήθελες εσύ;} \\
            }
        }
    }
    & \\ \\

    \theutterance \stepcounter{utterance}  
    & & & \multicolumn{4}{p{0.6\linewidth}}{
        \cellcolor[rgb]{0.9,0.9,0.9}{
            \makecell[{{p{\linewidth}}}]{
                \texttt{\tiny{[GM$\langle$P2]}}
                \texttt{{[}Πρόταση: 3 λουλούδια, 1 μπανάνα{]}} \\
            }
        }
    }
     \\ \\

    \theutterance \stepcounter{utterance}  
    & & & \multicolumn{2}{p{0.3\linewidth}}{
        \cellcolor[rgb]{0.9,0.9,0.9}{
            \makecell[{{p{\linewidth}}}]{
                \texttt{\tiny{[GM$|$GM]}}
                \texttt{proposal} \\
            }
        }
    }
    & & \\ \\

    \theutterance \stepcounter{utterance}  
    & & & \multicolumn{2}{p{0.3\linewidth}}{
        \cellcolor[rgb]{0.9,0.9,0.9}{
            \makecell[{{p{\linewidth}}}]{
                \texttt{\tiny{[GM$|$GM]}}
                [0, 3, 1]
            }
        }
    }
    & & \\ \\

    \theutterance \stepcounter{utterance}  
    & & \multicolumn{4}{p{0.6\linewidth}}{
        \cellcolor[rgb]{0.9,0.9,0.9}{
            \makecell[{{p{\linewidth}}}]{
                \texttt{\tiny{[P1$\langle$GM]}}
                \texttt{Ο άλλος παίκτης έχει υποβάλει μια μυστική πρόταση. Τώρα είναι η σειρά σας να υποβάλετε μια δική σας πρόταση.} \\
            }
        }
    }
    & \\ \\

    \theutterance \stepcounter{utterance}  
    & \multicolumn{4}{p{0.6\linewidth}}{
        \cellcolor[rgb]{0.9,0.9,0.9}{
            \makecell[{{p{\linewidth}}}]{
                \texttt{\tiny{[P1$\rangle$GM]}}
                \texttt{{[}Πρόταση: 2 λάμπες, 1 μπανάνα{]}} \\
            }
        }
    }
    & & \\ \\

    \theutterance \stepcounter{utterance}  
    & & & \multicolumn{2}{p{0.3\linewidth}}{
        \cellcolor[rgb]{0.9,0.9,0.9}{
            \makecell[{{p{\linewidth}}}]{
                \texttt{\tiny{[GM$|$GM]}}
                \texttt{proposal} \\
            }
        }
    }
    & & \\ \\

    \theutterance \stepcounter{utterance}  
    & & & \multicolumn{2}{p{0.3\linewidth}}{
        \cellcolor[rgb]{0.9,0.9,0.9}{
            \makecell[{{p{\linewidth}}}]{
                \texttt{\tiny{[GM$|$GM]}}
                [2, 0, 1]
            }
        }
    }
    & & \\ \\

    \theutterance \stepcounter{utterance}  
    & & & \multicolumn{2}{p{0.3\linewidth}}{
        \cellcolor[rgb]{0.9,0.9,0.9}{
            \makecell[{{p{\linewidth}}}]{
                \texttt{\tiny{[GM$|$GM]}}
                [[2, 0, 1], [0, 3, 1]]
            }
        }
    }
    & & \\ \\

\end{supertabular}
}

\end{document}
